\begin{definition}[Periodic checkerboard tiling]%
    \label{def:peps_kitaev_periodic_tiling}
    \lean{TNLean.PEPS.kitaevCornerEquiv,
        TNLean.PEPS.kitaevDoubleCoordinateEquiv,
        TNLean.PEPS.kitaevPeriodicTilingEquiv,
        TNLean.PEPS.kitaevTiledRight, TNLean.PEPS.kitaevTiledUp,
        TNLean.PEPS.kitaevTiledLeft, TNLean.PEPS.kitaevTiledDown,
        TNLean.PEPS.kitaevPeriodicTurned,
        TNLean.PEPS.kitaevPeriodicElementarySite}
    \leanok
    \uses{def:peps_kitaev_checkerboard_tensor}
    Let $w,h$ be positive integers. Divide the fine torus
    $\mathbb Z_{2w}\times\mathbb Z_{2h}$ into the disjoint squares
    based at $(2x,2y)$, with $(x,y)\in\mathbb Z_w\times\mathbb Z_h$.
    Number their corners clockwise from the upper right, so that
    $c_0=(1,1)$, $c_1=(1,0)$, $c_2=(0,0)$ and $c_3=(0,1)$.
    The site correspondence is $E(v,i)=2v+c_i$.
    Place $T_1$ at corners $0,2$ and $T_0$ at corners $1,3$.
    Contract every periodic bond with the identity and denote the resulting
    untwisted PEPS coefficient by $\Psi_T(\sigma)$.
    This is the periodic version of the elementary tensor in
    \cite[Section~7, ``Kitaev's code state'']{Schuch2010PEPS}.
\end{definition}

\begin{theorem}[The tiling respects periodic nearest neighbors]%
    \label{thm:peps_kitaev_periodic_tiling}
    \lean{TNLean.PEPS.kitaevDoubleCoordinateEquiv_apply,
        TNLean.PEPS.kitaevTiledRight_left, TNLean.PEPS.kitaevTiledUp_down,
        TNLean.PEPS.kitaevPeriodicTilingEquiv_right,
        TNLean.PEPS.kitaevPeriodicTilingEquiv_up,
        TNLean.PEPS.kitaevPeriodicTilingEquiv_left,
        TNLean.PEPS.kitaevPeriodicTilingEquiv_down,
        TNLean.PEPS.kitaevPeriodicTurned_right,
        TNLean.PEPS.kitaevPeriodicTurned_up}
    \leanok
    \uses{def:peps_kitaev_periodic_tiling}
    The correspondence $E$ is a bijection and identifies each horizontal
    or vertical step between tiled corners with the corresponding
    nearest-neighbor step on the fine torus. The orientation of the
    elementary tensor changes across every such bond, including bonds
    crossing a periodic seam.
\end{theorem}

\begin{proof}\leanok
    In each coordinate, division by two gives a unique coarse coordinate
    and a remainder in $\{0,1\}$. Adding one changes remainder $0$ to $1$;
    from remainder $1$ it advances the coarse coordinate and returns the
    remainder to $0$. The same rule applies at the periodic seam because
    the fine period is twice the coarse period. The two possible tensor
    orientations alternate in both cases.
\end{proof}

\begin{center}
    % Source: SCP10 Section 7.1, actual neighboring 2x2 checkerboard blocks.
    % All eight physical indices are fixed. Every tensor has virtual ports
    % top/right/bottom/left = 90/0/270/180. Ten internal bonds and twelve
    % open virtual ports give boundary signature (V,P)=(12,0).
    % The long horizontal bonds are the two labels crossing between blocks.
    \tenkzkernel
    \begin{tenkz}[rows={wire,wire},cols=5,bonds=none]
        \tn[at=(1,1),skin=box,name=A,ports={0:virtual,90:virtual,180:virtual,270:virtual}]{T_0}
        \tn[at=(1,2),skin=box,name=B,ports={0:virtual,90:virtual,180:virtual,270:virtual}]{T_1}
        \tn[at=(2,1),skin=box,name=C,ports={0:virtual,90:virtual,180:virtual,270:virtual}]{T_1}
        \tn[at=(2,2),skin=box,name=D,ports={0:virtual,90:virtual,180:virtual,270:virtual}]{T_0}
        \tn[at=(1,4),skin=box,name=E,ports={0:virtual,90:virtual,180:virtual,270:virtual}]{T_0}
        \tn[at=(1,5),skin=box,name=F,ports={0:virtual,90:virtual,180:virtual,270:virtual}]{T_1}
        \tn[at=(2,4),skin=box,name=G,ports={0:virtual,90:virtual,180:virtual,270:virtual}]{T_1}
        \tn[at=(2,5),skin=box,name=H,ports={0:virtual,90:virtual,180:virtual,270:virtual}]{T_0}
        \tnwire{A.0}{B.180}
        \tnwire{C.0}{D.180}
        \tnwire{A.270}{C.90}
        \tnwire{B.270}{D.90}
        \tnwire{E.0}{F.180}
        \tnwire{G.0}{H.180}
        \tnwire{E.270}{G.90}
        \tnwire{F.270}{H.90}
        \tnwire{B.0}{E.180}
        \tnwire{D.0}{G.180}
        \tnwire{open n}{A.90}
        \tnwire{open n}{B.90}
        \tnwire{open n}{E.90}
        \tnwire{open n}{F.90}
        \tnwire{C.270}{open s}
        \tnwire{D.270}{open s}
        \tnwire{G.270}{open s}
        \tnwire{H.270}{open s}
        \tnwire{open w}{A.180}
        \tnwire{open w}{C.180}
        \tnwire{F.0}{open e}
        \tnwire{H.0}{open e}
    \end{tenkz}
\end{center}
The eight physical indices are fixed. The two long bonds connect neighboring blocks. Their labels are read
from top to bottom on the eastern boundary of the left block and from
bottom to top on the western boundary of the right block. The vertical
pairs obey the corresponding reversed-order convention.

\begin{definition}[Periodic contraction of the four-site blocks]%
    \label{def:peps_kitaev_periodic_block_contraction}
    \lean{TNLean.PEPS.kitaevPeriodicBlockBoundary,
        TNLean.PEPS.kitaevPeriodicBlockedCoeff,
        TNLean.PEPS.kitaevPeriodicSiteLegs,
        TNLean.PEPS.kitaevPeriodicFineCoeff}
    \leanok
    \uses{def:peps_kitaev_checkerboard_block,def:peps_kitaev_periodic_tiling}
    Let $H_v,V_v\in\mathbb Z_2^2$ be the horizontal and vertical crossing
    pairs of a block $v$. Its clockwise boundary is
    $\alpha_v=(V_v,H_v,\operatorname{swap}V_{v-e_y},
        \operatorname{swap}H_{v-e_x})$.
    The periodic blocked coefficient is
    $\Psi_B(\tau)=\sum_{H,V}\prod_v B(\alpha_v;\tau_v)$.
    The fine coefficient $\Psi_T$ is the sum over these crossing pairs and
    the four independent internal bonds of every block, with one factor
    $T_0$ or $T_1$ for each of the four corners.
\end{definition}

\begin{theorem}[Global contraction of the elementary checkerboard PEPS]%
    \label{thm:peps_kitaev_global_blocking}
    \lean{TNLean.PEPS.kitaevPeriodicFineCoeff_eq_blocked,
        TNLean.PEPS.kitaevPeriodicFineCoeff_eq_colorNetwork,
        TNLean.PEPS.kitaevPeriodicBlockedCoeff_eq_colorNetwork,
        TNLean.PEPS.torusBondNetwork_kitaevPeriodicElementarySite_eq_fineCoeff,
        TNLean.PEPS.torusBondNetwork_kitaevPeriodicElementarySite_eq_colorNetwork,
        TNLean.PEPS.torusBondNetwork_kitaevPeriodicElementarySite_eq_quantumDoubleNetwork}
    \leanok
    \uses{def:peps_kitaev_periodic_tiling,
        thm:peps_kitaev_periodic_tiling,
        def:peps_kitaev_periodic_block_contraction,
        thm:peps_kitaev_checkerboard_blocking}
    For every fine physical configuration $\sigma$, put
    $\sigma_v=(\sigma(E(v,0)),\ldots,\sigma(E(v,3)))$.
    Let $K(t,r,d,l;\tau)=\delta_{\tau,F(t,r,d,l)}$ be the
    color-difference tensor. Then
    \begin{align}
        \Psi_T(\sigma)
         & =\sum_{a,b:\,\mathbb Z_w\times\mathbb Z_h\to\mathbb Z_2}
        \prod_v K(b_v,a_v,b_{v-e_y},a_{v-e_x};\sigma_v).
        \label{eq:peps_kitaev_global_blocking}
    \end{align}
    Thus the full fine-lattice state is the coarse color-difference PEPS
    under physical regrouping. The tensor $K$ is the dual toric-code tensor
    with physical order $(\sigma_3,\sigma_0,\sigma_1,\sigma_2)$.
    There is no additional scalar factor.
\end{theorem}

\begin{proof}\leanok
    The fine horizontal and vertical bonds partition into four internal
    bonds and four crossing bonds per block. Regroup the product by its
    four corners and distribute the independent internal sums to obtain
    one factor $B$ per coarse site. Write the horizontal crossing pair
    as $H_v$ and the vertical crossing pair as $V_v$. The clockwise
    boundary at $v$ is
    $(V_v,H_v,\operatorname{swap}V_{v-e_y},
        \operatorname{swap}H_{v-e_x})$.
    Theorem~\ref{thm:peps_kitaev_checkerboard_blocking} forces the two
    labels of every crossing pair to agree. Every pair with unequal
    labels makes the product zero. The surviving global bond assignments
    are therefore in bijection with $(a,b)$ through
    $H_v=(a_v,a_v)$ and $V_v=(b_v,b_v)$.
    Each contributes exactly the product in
    (\ref{eq:peps_kitaev_global_blocking}). The local color-difference
    identity identifies $K$ with the dual toric-code tensor.
\end{proof}

\begin{theorem}[Unitary physical regrouping of the global checkerboard state]%
    \label{thm:peps_kitaev_global_physical_blocking}
    \lean{TNLean.PEPS.kitaevPhysicalBlockingEquiv,
        TNLean.PEPS.kitaevPhysicalBlockingMatrix,
        TNLean.PEPS.kitaevPhysicalBlockingMatrix_isIsometry,
        TNLean.PEPS.kitaevPhysicalBlockingMatrix_conjTranspose_isIsometry,
        TNLean.PEPS.kitaevPhysicalBlockingMatrix_mulVec_checkerboard}
    \leanok
    \uses{thm:peps_kitaev_global_blocking}
    Let $R\ket{\sigma}=\ket{(\sigma_v)_v}$ be the physical regrouping
    from single fine spins to four-spin coarse sites. Then
    $R^\dagger R=1$, $RR^\dagger=1$ and $R\Psi_T=\Psi_K$.
\end{theorem}

\begin{proof}\leanok
    The disjoint tiling makes $\sigma\mapsto(\sigma_v)_v$ a bijection
    of computational bases, so both inner-product identities hold.
    Each coarse coefficient of $R\Psi_T$ has exactly one contributing
    fine configuration. Apply (\ref{eq:peps_kitaev_global_blocking}).
\end{proof}
% Scope: binary checkerboard tensors on even-period tori. This does not
% mark SCP10 Observation 6.6 for arbitrary regular G-isometric tensors complete.
